% ============================================================
% 3-р бүлэг: Судалгааны арга зүй
% ============================================================

\chapter{Судалгааны арга зүй}

% ======= 3.1 =======
\section{Судалгааны арга зүйн ерөнхий тойм}

Тус судалгааны ажил нь практик судалгаа (applied research) хэлбэрээр хэрэгжсэн бөгөөд Өвөрхангай аймгийн Хужирт сумын сувиллуудын одоогийн утсаар захиалга авах, цаасан бүртгэл хөтлөх үйл явцыг цахимжуулах зорилготой веб суурьтай мэдээллийн системийг хөгжүүлэхэд чиглэсэн. Хөгжүүлэлтийн аргачлалд Agile \cite{agilemanifesto2001} философийн зарчмуудыг баримтлан давтамжтай хөгжүүлэлтийн (iterative development) арга барилыг сонгосон. Энэхүү арга нь уламжлалт waterfall загвараас ялгаатай нь хэрэглэгчийн санал хүсэлтэд тулгуурлан системийг алхам алхмаар сайжруулах боломжийг олгодог.

Хөгжүүлэлтийн үйл явц нь долоон үндсэн үе шатаас бүрдсэн. Нэгдүгээр үе шатанд одоогийн сувиллын захиалгын процессын шинжилгээг хийж, утасаар болон биечлэн захиалга авах урсгалыг ажилтнуудтай ярилцлага хийх замаар нарийвчлан судалсан. Хоёрдугаар үе шатанд шаардлагын инженерчлэлийг гүйцэтгэж, функционал болон функционал бус шаардлагуудыг тодорхойлсон. Гуравдугаар үе шатанд системийн архитектурыг загварчилж, T3 Stack (Next.js, TypeScript, Prisma) технологийн стекийг сонгон full-stack веб системийн бүтцийг боловсруулсан. Дөрөвдүгээр үе шатанд өгөгдлийн сангийн схемийг Prisma ORM \cite{prisma2024} ашиглан тодорхойлж, PostgreSQL \cite{postgresql2024} дээр Supabase платформоор дамжуулан байршуулсан. Тавдугаар үе шатанд Next.js API Routes \cite{nextjs2024} болон NextAuth.js \cite{nextauth2024} ашиглан серверийн логик, нэвтрэлтийн системийг хөгжүүлсэн. Зургаадугаар үе шатанд React компонентууд, TailwindCSS ашиглан хэрэглэгчийн интерфейсийг бүтээсэн. Долоодугаар үе шатанд нэгтгэлийн туршилт, ачааллын туршилт, хэрэглэгчийн сэтгэл ханамжийн судалгааг гүйцэтгэсэн.

Давталт бүрийн хугацаа 2 долоо хоног байсан бөгөөд давталт бүрийн эцэст ажиллагаатай хувилбарыг (working increment) гаргаж, удирдагч багш болон сувиллын ажилтнуудаас санал хүсэлт авч дараагийн давталтын ажлыг төлөвлөсөн. Энэхүү давтамжтай арга нь хөгжүүлэлтийн явцад гарсан техникийн болон хэрэглэгчийн туршлагын асуудлуудыг эрт илрүүлж засварлах боломжийг олгосон.

% ======= 3.2 =======
\section{Судалгааны асуултууд}

Тус судалгааны ажлын хүрээнд дараах гурван судалгааны асуултыг дэвшүүлж, туршилтын үр дүнгээр хариулахыг зорьсон.

\textbf{RQ1 (Гүйцэтгэл):} Вэб суурьтай захиалгын систем нь одоогийн цаасан бүртгэл, утасаар захиалга авах процесстой харьцуулахад захиалга боловсруулах хугацааг статистик ач холбогдолтойгоор бууруулж чадах уу?

Энэхүү асуулт нь системийн гол зорилго болох үйл ажиллагааны үр ашгийг хэмжихэд чиглэсэн. Хариултыг тодорхойлохдоо нэг захиалга боловсруулах, мэдээлэл хайх, тайлан гаргах зэрэг үйлдлүүдийн хугацааг хуучин болон шинэ системд хэмжиж харьцуулсан.

\textbf{RQ2 (Алдаа):} Систем нь давхар захиалга (double-booking), огнооны зөрөө зэрэг хүний алдаанаас үүдэлтэй асуудлуудыг бүрэн арилгаж чадах уу?

Цаасан бүртгэлийн хамгийн том сул тал нь ижил өрөөг ижил хугацаанд хоёр өөр хэрэглэгчид захиалахад системийн хяналт байхгүйгээс давхардал үүсэх явдал юм. Энэ асуултад хариулахдаа зэрэг хандалтын (concurrency) туршилтаар системийн pessimistic locking механизмыг шалгасан.

\textbf{RQ3 (Хэрэглэгчийн сэтгэл ханамж):} Системийн ашиглах чадварын (usability) үнэлгээ нь System Usability Scale (SUS) \cite{brooke1996sus} аргачлалын хүлээн зөвшөөрөгдөхүйц босго болох 68 онооноос дээш үнэлгээ авч чадах уу?

SUS нь олон улсад өргөн хэрэглэгддэг хэрэглэгчийн туршлагын (UX) хэмжүүр бөгөөд 10 асуултын 5 оноотой Likert scale-ээр үнэлгээ авч 0--100 хүртэлх оноогоор илэрхийлдэг. Bangor нарын судалгаагаар \cite{bangor2008sus} 68-аас дээш оноо нь ``хүлээн зөвшөөрөгдөхүйц'' (acceptable) гэж тооцогддог.

% ======= 3.3 =======
\section{Системийн үйл ажиллагааны тухай дэлгэрэнгүй}

\subsection{Системд ерөнхий бүтэц}

Энэхүү систем нь T3 Stack дээр суурилсан бүрэн хэмжээний веб систем бөгөөд Next.js, Prisma ORM, NextAuth.js технологиудыг ашигласан. Хэрэглэгчийн хүсэлт нь Next.js App Router-аар дамжин Prisma ORM-р PostgreSQL өгөгдлийн санд хандах бүтэцтэй ба T3 технологийн стек нь TypeScript хэлийг ашиглан frontend, backend хооронд кодын давхардлыг багасгаж, системийн найдвартай байдлыг хангадаг.

\begin{figure}[htbp]
\centering
\begin{tikzpicture}[
    node distance=0.9cm,
    font=\normalsize,
    layer/.style={
        rectangle,
        rounded corners=6pt,
        draw,
        line width=1.2pt,
        minimum width=11cm,
        minimum height=1.6cm,
        align=center,
        inner sep=6pt
    },
    client/.style   ={layer, draw=violet!60,      fill=violet!12},
    frontend/.style ={layer, draw=orange!70,      fill=orange!15},
    backend/.style  ={layer, draw=red!55,         fill=red!12},
    database/.style ={layer, draw=green!55!black, fill=green!15},
    arr/.style={-{Latex[length=3mm, width=2mm]}, line width=1.2pt}
]
\node[client] (client)
    {\textbf{Хэрэглэгчийн давхарга (client)} \\[2pt]
     Зочин \ $\cdot$ \ Админ \ $\cdot$ \ Ресепшн};

\node[frontend, below=of client] (frontend)
    {\textbf{Харагдах хэсэг (frontend)} \\[2pt]
     Next.js 14 \ $\cdot$ \ React \ $\cdot$ \ TailwindCSS};

\node[backend, below=of frontend] (backend)
    {\textbf{Серверийн хэсэг (backend)} \\[2pt]
     Next.js API Routes \ $\cdot$ \ NextAuth.js \ $\cdot$ \ Prisma ORM};

\node[database, below=of backend] (database)
    {\textbf{Өгөгдлийн давхарга (database)} \\[2pt]
     PostgreSQL \ $\cdot$ \ Supabase};

\draw[arr, violet!60] (client)   -- (frontend);
\draw[arr, orange!70] (frontend) -- (backend);
\draw[arr, red!55]    (backend)  -- (database);
\end{tikzpicture}
\caption{Системийн дөрвөн давхаргат бүтэц}
\label{fig:architecture}
\end{figure}

\subsection{Системийн үндсэн урсгал}

Системийн үндсэн үйл ажиллагааны урсгал нь дараах дарааллаар явагдана: зочин буюу үйлчлүүлэгч вэб хуудсанд орж сувиллын мэдээлэл болон өрөөний жагсаалттай танилцана, захиалга хийхийн тулд ирэх унаа, огноо болон хүний тоо сонгон боломжтой өрөөг хайлт хийж, захиалга үүсгэнэ. Үүсгэсэн захиалга дээр амжилттай төлбөр төлөгдсөн тохиолдолд систем хэрэглэгчид мэдэгдэл илгээж захиалга бүртгэгдсэн гэдгийг баталгаажуулна. Энэхүү урсгал нь одоогийн утсаар захиалга авах, цаасан бүртгэл хөтлөх уламжлалт аргыг цахим орчинд шилжүүлж байгаа юм.

\begin{figure}[H]
  \centering
  \figplaceholder[0.55]{system_flow}
  \caption{Системийн үндсэн урсгалын диаграмм}
  \label{fig:system-flow}
\end{figure}

% ======= 3.4 =======
\section{Системийн хэрэглэгчид}

Системд дөрвөн үүргийн (role) хэрэглэгч байх бөгөөд үүрэг бүр нь тодорхой эрх, боломжоор хязгаарлагдана. Зочин (Guest) хэрэглэгч бүртгэлгүй байдлаар сувиллын ерөнхий мэдээлэл, өрөөний жагсаалт, үнийн мэдээллийг харж болох боловч захиалга хийхийн тулд бүртгүүлэх шаардлагатай болно. Энэхүү нээлттэй танилцах боломж нь боломжит үйлчлүүлэгчдийг татах маркетингийн чиг үүрэгтэй.

Үйлчлүүлэгч гэдэг нь бүртгүүлсний дараа идэвхждэг үүрэг бөгөөд захиалга үүсгэх, өөрийн захиалгын түүх харах, профайлын мэдээлэл засах, баталгаажсан захиалгаас хугацаанаас өмнө цуцлах боломжтой. Ажилтан нь өөрийн сувиллын орж ирсэн захиалгуудыг хянаж баталгаажуулах эсвэл цуцлах, ирэх үйлчлүүлэгчдийн нэрсийн жагсаалт гаргах үүрэгтэй. Захиалгын цаг хугацаа, нийт орны хүчин чадлыг ажилтан хянаснаар давхардсан захиалга болон алдааг урьдчилан сэргийлнэ. Админ нь системийн хамгийн өндөр эрхтэй үүрэг бөгөөд бүх сувиллын захиалга, хэрэглэгч, өрөөний мэдээлэл удирдах, тайлан статистик харах, ажилтан бүртгэх боломжтой.

\vspace{0.5cm}

\begin{table}[H]
  \centering
  \caption{Хэрэглэгчийн эрхийн хүснэгт}
  \label{tab:role-matrix}
  \begin{tabular}{|l|c|c|c|c|}
    \hline
    \textbf{Үйлдэл} & \textbf{Зочин} & \textbf{Үйлчлүүлэгч} & \textbf{Ажилтан} & \textbf{Админ} \\
    \hline
    Сувиллын мэдээлэл харах  & + & + & + & + \\
    \hline
    Захиалга хийх            & -- & + & -- & + \\
    \hline
    Өөрийн захиалга цуцлах   & -- & + & -- & + \\
    \hline
    Захиалга баталгаажуулах   & -- & -- & + & + \\
    \hline
    Ажилтан бүртгэх          & -- & -- & -- & + \\
    \hline
    Тайлан харах             & -- & -- & -- & + \\
    \hline
  \end{tabular}
\end{table}

% ======= 3.5 =======
\section{Системийн шаардлага}

Судалгааны явцад системд шаардлагатай үндсэн функцуудыг сувиллын ажилтантай хийсэн асуулга, ярилцлагын аргаар тодорхойлсон.
Мөн системийн найдвартай байдал, аюулгүй байдлын шаардлагуудыг тусгайлан авч үзсэн. Эдгээр шаардлагуудыг тодорхойлохдоо Эрүүл мэндийн яамны нийтийн эрүүл мэнд, нөхөн сэргээх үйлчилгээний статистик мэдээллийг \cite{mohstats2023} суурь мэдээлэл болгон ашигласан.
Доорх зурагт Хужирт сумын нэгэн сувилалд 2016–2023 оны хооронд үйлчлүүлсэн үйлчлүүлэгчдийн тооны өөрчлөлтийг харуулав.
\begin{figure}[H]
  \centering
  \begin{tikzpicture}
    \begin{axis}[
      ybar,
      bar width=0.58cm,
      width=0.85\textwidth,
      height=6.5cm,
      xtick={1,2,3,4,5,6,7,8},
      xticklabels={2016,2017,2018,2019,2020,2021,2022,{2023/11}},
      nodes near coords,
      nodes near coords align={vertical},
      every node near coord/.style={font=\scriptsize},
      ymin=0,
      ymax=6500,
      ylabel={Үйлчлүүлэгчийн тоо (хүн)},
      xlabel={Он},
      enlarge x limits=0.08,
      tick label style={font=\small},
      label style={font=\small},
    ]
      \addplot[fill=blue!45, draw=blue!70] coordinates {
        (1,1450)(2,2937)(3,3346)(4,4069)
        (5,3863)(6,2521)(7,5146)(8,4247)
      };
    \end{axis}
  \end{tikzpicture}
  \caption{Хужирт сумын нэг сувиллын үйлчлүүлэгчийн тооны өөрчлөлт (2016--2023)\\[2pt]}
  \label{fig:customers-stats}
\end{figure}
\vspace{0.5cm}

\begin{table}[H]
  \centering
  \caption{Системийн шаардлагын жагсаалт}
  \label{tab:functional-req}
  \begin{tabular}{|l|l|p{6cm}|l|}
    \hline
    \textbf{FR-ID} & \textbf{Шаардлагын нэр} & \textbf{Тодорхойлолт} & \textbf{Хамаарах үүрэг} \\
    \hline
    FR-01 & Нэвтрэх               & Имэйл/нууц үгээр нэвтрэх                          & Бүгд \\
    \hline
    FR-02 & Бүртгүүлэх            & Шинэ хэрэглэгч үүсгэх, имэйл баталгаажуулалт                  & Зочин \\
    \hline
    FR-03 & Захиалга үүсгэх       & Өрөө, огноо, хүний тоо сонгон захиалга өгөх                   & Үйлчлүүлэгч \\
    \hline
    FR-04 & Захиалга баталгаажуулах & Ажилтан захиалгыг хянаж батлах                               & Ажилтан \\
    \hline
    FR-05 & Захиалга цуцлах       & Үйлчлүүлэгч эсвэл ажилтан захиалга цуцлах                    & Үйлчлүүлэгч, Ажилтан \\
    \hline
    FR-06 & Өрөөний хүртээмж      & Сонгосон огноонд хоосон өрөө шүүх                             & Зочин, Үйлчлүүлэгч \\
    \hline
    FR-07 & Тайлан гаргах         & Сарын захиалга, орлого, дүүргэлтийн хувь                      & Админ \\
    \hline
    FR-08 & Хэрэглэгч удирдах     & Хэрэглэгч, ажилтан бүртгэх, засах                             & Админ \\
    \hline
    FR-09 & Мэдэгдэл илгээх       & Захиалга баталгаажсан, ирэх өдрийг сануулах                   & Систем \\
    \hline
    FR-10 & Профайл засах         & Нэр, утас, нууц үг шинэчлэх                                   & Үйлчлүүлэгч \\
    \hline
  \end{tabular}
\end{table}

% ======= 3.6 =======

Статистик мэдээллээс харахад 2022 онд үйлчлүүлэгчдийн тоо 5,146-д хүрч, судлагдсан хугацаанд хамгийн өндөр үзүүлэлтэд хүрсэн байна. Энэхүү өсөлтийн хандлага нь системд ирэх ачаалал нэмэгдэхийг илтгэж байгаа тул уг өгөгдөлд тулгуурлан функционал бус шаардлагуудыг тодорхойлсон.
Нэгдүгээрт, жилийн дундаж болон оргил ачааллын үед нэгэн зэрэг нэвтрэх хэрэглэгчдийн тоо мэдэгдэхүйц өсөх хандлагатай тул \textbf{NFR-01 гүйцэтгэлийн шаардлага}-ын хүрээнд 100-аас дээш хэрэглэгчийг нэгэн зэрэг 200 ms-ээс бага хариу хугацаатайгаар дэмжих шаардлагыг тогтоосон.
Хоёрдугаарт, олон хэрэглэгч нэгэн зэрэг ижил өрөөг захиалах нөхцөл үүсэх магадлалтай учир \textbf{NFR-03 захиалга түгжих} механизмыг хэрэгжүүлэх шаардлагатай гэж үзсэн.
Гуравдугаарт, хэрэглэгчдийн тодорхой хувь нь гар утас болон таблет төхөөрөмжөөс системд хандах хэрэгцээтэй тул \textbf{NFR-04 responsive дизайн}-ыг зайлшгүй хангахаар тусгасан.
\vspace{0.5cm}
\begin{table}[H]
  \centering
  \caption{Функционал бус шаардлагууд}
  \label{tab:nonfunctional-req}
  \begin{tabular}{|l|l|p{5.5cm}|p{3.5cm}|}
    \hline
    \textbf{NFR-ID} & \textbf{Ангилал} & \textbf{Шаардлага} & \textbf{Хэмжих үзүүлэлт} \\
    \hline
    NFR-01 & Гүйцэтгэл      
    & API хариу хугацаа дунджаар 200 ms-ээс бага, 100+ хэрэглэгчийг нэгэн зэрэг дэмжих    
    & Ачааллын тест (Load test) \\
    \hline
    NFR-02 & Аюулгүй байдал 
    & HTTPS протокол, JWT токен, bcrypt алгоритмаар нууц үг хэшлэх, SQL injection хамгаалалт                  
    & Аюулгүй байдлын тест \\
    \hline
    NFR-03 & Өгөгдлийн бүрэн бүтэн байдал 
    & Захиалгыг баталгаажуулах хүртэл 10 минутын хугацаанд түгжих механизм хэрэгжүүлэх                            
    & Зэрэг хандалтын тест \\
    \hline
    NFR-04 & Хүртээмж       
    & Responsive дизайн (гар утас, таблет, компьютер дээр зохицох)          
    & Дэлгэцийн тест \\
    \hline
    NFR-05 & Засвар үйлчилгээ 
    & Лог бүртгэл, кодын баримтжуулалт, хялбар deploy хийх боломж                              
    & Кодын хяналт \\
    \hline
  \end{tabular}
\end{table}

% ======= 3.7 =======
\section{Юзкейс диаграммууд}

\subsection{Үйлчлүүлэгчийн юзкейс диаграмм}

Үйлчлүүлэгч бүртгүүлсний дараа захиалга үүсгэх, өөрийн захиалгын түүх болон статусыг харах, профайлын мэдээлэл засах, хугацааны нөхцөлийг хангасан захиалгаа цуцлах боломжтой.

\begin{figure}[H]
\centering
\begin{tikzpicture}[
    font=\normalsize,
    usecase/.style={
        ellipse,
        draw=blue!55!black,
        line width=0.9pt,
        fill=blue!18,
        minimum width=4.2cm,
        minimum height=1.6cm,
        align=center,
        inner sep=3pt
    },
    arr/.style={-{Latex[length=3mm, width=2mm]}, line width=0.9pt},
    ext/.style={-{Latex[length=2.5mm, width=1.8mm]}, line width=0.7pt, dashed}
]

\draw[line width=0.9pt] (0, 0) circle (0.35cm);
\draw[line width=0.9pt] (0, -0.35) -- (0, -1.5);
\draw[line width=0.9pt] (-0.7, -0.85) -- (0.7, -0.85);
\draw[line width=0.9pt] (0, -1.5) -- (-0.5, -2.3);
\draw[line width=0.9pt] (0, -1.5) -- (0.5, -2.3);
\node[below=0.1cm] at (0, -2.4) {\textbf{Үйлчлүүлэгч}};

\node[usecase] (register) at (7.5,  4.8)   {Бүртгэл үүсгэх / \\ Нэвтрэх};
\node[usecase] (viewroom) at (7.5,  2.7)   {Өрөө харах / \\ Хайх};
\node[usecase] (book)     at (7.5,  0.7)   {Захиалга үүсгэх};
\node[usecase] (mybook)   at (7.5, -1.3)   {Захиалгаа харах / \\ Засах};
\node[usecase] (cancel)   at (7.5, -3.3)   {Захиалгаа цуцлах};
\node[usecase] (payment)  at (7.5, -5.3)   {Урьдчилгаа төлбөр \\ төлөх};

\draw[arr] (0.7, -0.85) -- (register.west);
\draw[arr] (0.7, -0.85) -- (viewroom.west);
\draw[arr] (0.7, -0.85) -- (book.west);
\draw[arr] (0.7, -0.85) -- (mybook.west);
\draw[arr] (0.7, -0.85) -- (cancel.west);
\draw[arr] (0.7, -0.85) -- (payment.west);

\draw[ext] (viewroom.south) -- (book.north);
\draw[ext] (payment.north east) to[bend right=30] (book.south east);

\end{tikzpicture}
\caption{Үйлчлүүлэгчийн юзкейс диаграмм}
\label{fig:usecase-customer}
\end{figure}

\subsection{Ажилтан болон Админы юзкейс диаграмм}

Ажилтан өөрийн сувиллд ирсэн захиалгуудыг хянаж баталгаажуулах, шаардлагатай тохиолдолд цуцлах, ирэх хугацааны үйлчлүүлэгчдийн жагсаалт гаргах боломжтой бөгөөд өрөөний хүртээмжийн хяналтыг удирддаг. Админ системийн хамгийн өндөр эрхийн хэрэглэгч бөгөөд бүх сувиллын захиалга, өрөөний мэдээлэл, хэрэглэгчийн бүртгэлийг удирдах, ажилтан нэмэх, тайлан статистик харах зэрэг бүх үйлдлийг гүйцэтгэх боломжтой.

\begin{figure}[H]
\centering
\begin{tikzpicture}[
    font=\normalsize,
    shared/.style={
        ellipse,
        draw=green!50!black,
        line width=0.9pt,
        fill=green!20,
        minimum width=4.2cm,
        minimum height=1.5cm,
        align=center,
        inner sep=3pt
    },
    adminuc/.style={
        ellipse,
        draw=orange!70!black,
        line width=0.9pt,
        fill=orange!25,
        minimum width=4.2cm,
        minimum height=1.5cm,
        align=center,
        inner sep=3pt
    },
    arr/.style={-{Latex[length=3mm, width=2mm]}, line width=0.9pt}
]

\draw[line width=0.9pt] (0,  4.5) circle (0.35cm);
\draw[line width=0.9pt] (0,  4.15) -- (0,  3.0);
\draw[line width=0.9pt] (-0.7, 3.65) -- (0.7, 3.65);
\draw[line width=0.9pt] (0,  3.0) -- (-0.5, 2.2);
\draw[line width=0.9pt] (0,  3.0) -- ( 0.5, 2.2);
\node[below=0.1cm, align=center] at (0, 2.1) {Хүлээн авах \\ ажилтан};

\draw[line width=0.9pt] (0, -1.5) circle (0.35cm);
\draw[line width=0.9pt] (0, -1.85) -- (0, -3.0);
\draw[line width=0.9pt] (-0.7, -2.35) -- (0.7, -2.35);
\draw[line width=0.9pt] (0, -3.0) -- (-0.5, -3.8);
\draw[line width=0.9pt] (0, -3.0) -- ( 0.5, -3.8);
\node[below=0.1cm] at (0, -3.9) {Админ};

\node[shared]  (confirm)  at (8,  5.2)   {Захиалга баталгаажуулах};
\node[shared]  (viewlist) at (8,  3.0)   {Захиалгын жагсаалт \\ харах};
\node[shared]  (status)   at (8,  0.8)   {Захиалгын төлөв \\ өөрчлөх};

\node[adminuc] (rooms)    at (8, -1.4)   {Өрөө / Үйлчилгээ \\ нэмэх, засах};
\node[adminuc] (staff)    at (8, -3.3)   {Ажилтан удирдах};
\node[adminuc] (report)   at (8, -5.2)   {Тайлан, статистик \\ гаргах};

\draw[arr] (0.7,  3.65) -- (confirm.west);
\draw[arr] (0.7,  3.65) -- (viewlist.west);
\draw[arr] (0.7,  3.65) -- (status.west);

\draw[arr] (0.7, -2.35) -- (status.west);
\draw[arr] (0.7, -2.35) -- (viewlist.west);
\draw[arr] (0.7, -2.35) -- (rooms.west);
\draw[arr] (0.7, -2.35) -- (staff.west);
\draw[arr] (0.7, -2.35) -- (report.west);

\end{tikzpicture}
\caption{Ажилтан болон Админы юзкейс диаграмм}
\label{fig:usecase-staff-admin}
\end{figure}

% ======= 3.8 =======
\section{Дарааллын диаграммууд}

\subsection{Бүртгэл ба нэвтрэлтийн дараалал}

Хэрэглэгч имэйл болон нууц үгээ оруулахад NextAuth.js хүсэлтийг хүлээн авч Prisma ORM-р дамжуулан PostgreSQL өгөгдлийн санд хадгалагдсан хэрэглэгчийн мэдээлэлтэй шалгана, амжилттай баталгаажсан тохиолдолд JWT session токен үүсгэн хэрэглэгчийн хөтөч дээр хадгалагдана.

\begin{figure}[H]
\centering
\resizebox{0.95\textwidth}{!}{
\begin{tikzpicture}[
    font=\small,
    actor/.style={
        rectangle,
        draw=blue!40!black,
        fill=blue!8,
        rounded corners=4pt,
        minimum width=2.6cm,
        minimum height=0.7cm,
        align=center
    },
    line/.style={dashed, gray!80, line width=0.6pt},
    msg/.style={-Latex, line width=0.8pt},
    ret/.style={-Latex, dashed, line width=0.8pt}
]

\def\xa{0}
\def\xb{4.3}
\def\xc{8.6}
\def\xd{12.9}

% Actors
\node[actor] (u)  at (\xa, 0) {Хэрэглэгч \\ (веб)};
\node[actor] (na) at (\xb, 0) {NextAuth.js};
\node[actor] (pr) at (\xc, 0) {Prisma ORM};
\node[actor] (db) at (\xd, 0) {PostgreSQL};

% Lifelines
\draw[line] (\xa, -0.6) -- (\xa, -13.2);
\draw[line] (\xb, -0.6) -- (\xb, -13.2);
\draw[line] (\xc, -0.6) -- (\xc, -13.2);
\draw[line] (\xd, -0.6) -- (\xd, -13.2);

% [БҮРТГЭЛ]
\node[anchor=west, font=\bfseries\small, text=gray!80!black] at (-0.5, -1.0) {Бүртгэл};

\draw[msg] (\xa, -1.8) -- node[above, font=\scriptsize] {1. Бүртгэлийн мэдээлэл (нэр, имэйл, нууц үг)} (\xb, -1.8);
\draw[msg] (\xb, -2.6) -- node[above, font=\scriptsize] {2. Шинэ хэрэглэгч үүсгэх хүсэлт} (\xc, -2.6);
\draw[msg] (\xc, -3.4) -- node[above, font=\scriptsize] {3. INSERT INTO users} (\xd, -3.4);

\draw[ret] (\xd, -4.2) -- node[above, font=\scriptsize] {4. Хэрэглэгч амжилттай үүссэн} (\xc, -4.2);
\draw[ret] (\xc, -5.0) -- node[above, font=\scriptsize] {5. User объект буцаах} (\xb, -5.0);
\draw[ret] (\xb, -5.8) -- node[above, font=\scriptsize] {6. Бүртгэл амжилттай (нэвтрэх хуудас руу)} (\xa, -5.8);

% Separator
\draw[loosely dashed, gray!50] (-0.8, -6.6) -- (\xd+0.8, -6.6);

% [НЭВТРЭЛТ]
\node[anchor=west, font=\bfseries\small, text=gray!80!black] at (-0.5, -7.2) {Нэвтрэлт};

\draw[msg] (\xa, -8.0) -- node[above, font=\scriptsize] {7. Имэйл, нууц үг оруулах} (\xb, -8.0);
\draw[msg] (\xb, -8.8) -- node[above, font=\scriptsize] {8. Имэйлээр хэрэглэгч хайх} (\xc, -8.8);
\draw[msg] (\xc, -9.6) -- node[above, font=\scriptsize] {9. SELECT * FROM users WHERE email=?} (\xd, -9.6);

\draw[ret] (\xd, -10.4) -- node[above, font=\scriptsize] {10. Хэрэглэгчийн мэдээлэл} (\xc, -10.4);
\draw[ret] (\xc, -11.2) -- node[above, font=\scriptsize] {11. User + hashed нууц үг} (\xb, -11.2);

\node[draw=gray!50, fill=yellow!15, text=black, align=center, font=\scriptsize, rounded corners=3pt, inner sep=4pt] at (\xb, -12.0) {bcrypt.compare() \\ нууц үг шалгах};

\draw[ret] (\xb, -12.8) -- node[above, font=\scriptsize] {12. JWT session токен (cookie-д хадгалах)} (\xa, -12.8);

\end{tikzpicture}
}
\caption{Бүртгэл ба нэвтрэлтийн дараалал диаграмм}
\label{fig:seq-auth}
\end{figure}

\subsection{Захиалга үүсгэх дараалал}

Үйлчлүүлэгч огноо, өрөө, хүний тоо сонгон захиалга илгээхэд tRPC mutation дуудагдаж серверийн хэсэгт хүсэлт очно, Prisma ашиглан өгөгдлийн санд шинэ захиалгын бичлэг үүсч, захиалгын статус ``хүлээгдэж байна'' (pending) болж тохируулагдана.

\begin{figure}[H]
\centering
\resizebox{0.95\textwidth}{!}{
\begin{tikzpicture}[
    font=\small,
    actor/.style={
        rectangle,
        draw=blue!40!black,
        fill=blue!8,
        rounded corners=4pt,
        minimum width=2.6cm,
        minimum height=0.7cm,
        align=center
    },
    line/.style={dashed, gray!80, line width=0.6pt},
    msg/.style={-Latex, line width=0.8pt},
    ret/.style={-Latex, dashed, line width=0.8pt}
]

\def\xa{0}
\def\xb{4.8}
\def\xc{9.6}
\def\xd{14.4}

% Actors
\node[actor] (u)  at (\xa, 0) {Үйлчлүүлэгч \\ (веб)};
\node[actor] (na) at (\xb, 0) {tRPC \\ (API Router)};
\node[actor] (pr) at (\xc, 0) {Prisma ORM};
\node[actor] (db) at (\xd, 0) {PostgreSQL};

% Lifelines
\draw[line] (\xa, -0.6) -- (\xa, -8.5);
\draw[line] (\xb, -0.6) -- (\xb, -8.5);
\draw[line] (\xc, -0.6) -- (\xc, -8.5);
\draw[line] (\xd, -0.6) -- (\xd, -8.5);

\draw[msg] (\xa, -1.3) -- node[above, font=\scriptsize] {1. Огноо, өрөө, хүний тоо сонгон захиалга илгээх} (\xb, -1.3);

\draw[msg] (\xb, -2.9) -- node[above, font=\scriptsize] {2. booking.create mutation дуудах} (\xc, -2.9);
\draw[msg] (\xc, -3.7) -- node[above, font=\scriptsize] {3. Өрөөний хүртээмж шалгах (SELECT)} (\xd, -3.7);

\draw[ret] (\xd, -4.5) -- node[above, font=\scriptsize] {4. Өрөө хоосон байна} (\xc, -4.5);

\draw[msg] (\xc, -5.3) -- node[above, font=\scriptsize] {5. INSERT INTO bookings (status=pending)} (\xd, -5.3);

\draw[ret] (\xd, -6.1) -- node[above, font=\scriptsize] {6. Захиалга амжилттай үүссэн} (\xc, -6.1);
\draw[ret] (\xc, -6.9) -- node[above, font=\scriptsize] {7. Booking объект буцаах} (\xb, -6.9);
\draw[ret] (\xb, -7.7) -- node[above, font=\scriptsize] {8. Захиалга үүссэн (status: pending)} (\xa, -7.7);

\end{tikzpicture}
}
\caption{Захиалга үүсгэх дараалал диаграмм}
\label{fig:seq-booking}
\end{figure}

\subsection{Захиалга баталгаажуулах дараалал}

Ажилтан хүлээгдэж буй захиалгуудын жагсаалтаас захиалга сонгон баталгаажуулахад tRPC mutation дуудагдаж өгөгдлийн санд захиалгын статус ``баталгаажсан'' (confirmed) болж шинэчлэгдэх бөгөөд систем автоматаар үйлчлүүлэгчид мэдэгдэл илгээнэ.

\begin{figure}[H]
\centering
\resizebox{0.95\textwidth}{!}{
\begin{tikzpicture}[
    font=\small,
    actor/.style={
        rectangle,
        draw=blue!40!black,
        fill=blue!8,
        rounded corners=4pt,
        minimum width=2.6cm,
        minimum height=0.7cm,
        align=center
    },
    line/.style={dashed, gray!80, line width=0.6pt},
    msg/.style={-Latex, line width=0.8pt},
    ret/.style={-Latex, dashed, line width=0.8pt}
]

\def\xa{0}
\def\xb{4.8}
\def\xc{9.6}
\def\xd{14.4}

% Actors
\node[actor] (st)  at (\xa, 0) {Ажилтан \\ (Дэлгэц)};
\node[actor] (na) at (\xb, 0) {tRPC \\ (API Router)};
\node[actor] (pr) at (\xc, 0) {Prisma ORM};
\node[actor] (db) at (\xd, 0) {PostgreSQL};

% Lifelines
\draw[line] (\xa, -0.6) -- (\xa, -8.3);
\draw[line] (\xb, -0.6) -- (\xb, -8.3);
\draw[line] (\xc, -0.6) -- (\xc, -8.3);
\draw[line] (\xd, -0.6) -- (\xd, -8.3);

\draw[msg] (\xa, -1.3) -- node[above, font=\scriptsize] {1. Захиалга баталгаажуулах хүсэлт илгээх} (\xb, -1.3);
\draw[msg] (\xb, -2.3) -- node[above, font=\scriptsize] {2. update.booking(status='confirmed')} (\xc, -2.3);
\draw[msg] (\xc, -3.3) -- node[above, font=\scriptsize] {3. UPDATE bookings SET status='confirmed'} (\xd, -3.3);

\draw[ret] (\xd, -4.3) -- node[above, font=\scriptsize] {4. Бичлэг амжилттай шинэчлэгдсэн} (\xc, -4.3);
\draw[ret] (\xc, -5.3) -- node[above, font=\scriptsize] {5. Захиалгын объект буцаах} (\xb, -5.3);

\node[draw=gray!50, fill=yellow!15, text=black, align=center, font=\scriptsize, rounded corners=3pt, inner sep=4pt] at (\xb, -6.3) {Үйлчлүүлэгч рүү баталгаажсан \\ тухай автомат имэйл илгээх};

\draw[ret] (\xb, -7.3) -- node[above, font=\scriptsize] {6. Дэлгэцэд "Амжилттай" төлөв харуулах} (\xa, -7.3);

\end{tikzpicture}
}
\caption{Захиалга баталгаажуулах дараалал диаграмм}
\label{fig:seq-confirm}
\end{figure}

% ======= 3.9 =======
\section{Өгөгдлийн сангийн ерөнхий загвар}

\subsection{Үндсэн хүснэгтүүд ба харилцаа холбоо}

Өгөгдлийн сангийн бүтэц нь Prisma schema дээр суурилсан таван үндсэн хүснэгтээс бүрдэх бөгөөд тэдгээрийн хоорондын харилцаа нь гадаад түлхүүр (foreign key) ашиглан нэг-олон (one-to-many) хамаарлаар холбогдоно. Хүснэгтүүд нь хэрэглэгч, сувилл, өрөө, захиалга, эмчилгээний бүртгэлийн мэдээллийг тусад нь хадгалж системийн мэдээллийн бүрэн бүтэн байдлыг хангана.

\begin{figure}[H]
\centering
\resizebox{\textwidth}{!}{
\begin{tikzpicture}[
    font=\small,
    % ---- Entity (хүснэгт) хэв маяг ----
    entity/.style={
        rectangle split,
        rectangle split parts=2,
        rectangle split part fill={violet!25, white},
        draw=violet!40,
        line width=0.8pt,
        rounded corners=3pt,
        align=left,
        text width=4.8cm,
        inner sep=6pt
    },
    % ---- Харилцааны шугам ----
    rel/.style={line width=0.7pt, black!70},
    % ---- Multiplicity label ----
    mult/.style={
        font=\small,
        inner sep=3pt,
        fill=white,
        draw=none
    }
]

% ---- Талбар харуулах макронууд ----
\newcommand{\pk}[2]{\underline{#1} : \textit{#2} \ \textsc{\footnotesize(PK)} \\}
\newcommand{\fk}[2]{#1 : \textit{#2} \ \textsc{\footnotesize(FK)} \\}
\newcommand{\attr}[2]{#1 : \textit{#2} \\}

% ---------- ROW 1: users, sanatoriums, rooms ----------
\node[entity] (User) at (0, 10) {
    \textbf{users} \nodepart{second}
    \pk{id}{String}
    \attr{name}{String}
    \attr{email}{String (unique)}
    \attr{password}{String}
    \attr{role}{Role}
    \attr{phone}{String?}
    \attr{createdAt}{DateTime}
};

\node[entity] (Sana) at (6.5, 10) {
    \textbf{sanatoriums} \nodepart{second}
    \pk{id}{String}
    \attr{name}{String}
    \attr{location}{String}
    \attr{description}{String?}
    \attr{phone}{String?}
    \attr{createdAt}{DateTime}
};

\node[entity] (Room) at (13, 10) {
    \textbf{rooms} \nodepart{second}
    \pk{id}{String}
    \fk{sanatoriumId}{String}
    \attr{name}{String}
    \attr{type}{RoomType}
    \attr{capacity}{Int}
    \attr{pricePerNight}{Float}
    \attr{imageUrl}{String?}
};

% ---------- ROW 2: bookings, booking_rooms ----------
\node[entity] (Booking) at (0, 2) {
    \textbf{bookings} \nodepart{second}
    \pk{id}{String}
    \fk{userId}{String}
    \attr{checkIn}{DateTime}
    \attr{checkOut}{DateTime}
    \attr{guests}{Int}
    \attr{totalPrice}{Float}
    \attr{status}{Status}
    \attr{createdAt}{DateTime}
};

\node[entity] (BR) at (13, 2) {
    \textbf{booking\_rooms} \nodepart{second}
    \pk{id}{String}
    \fk{bookingId}{String}
    \fk{roomId}{String}
    \attr{pricePerNight}{Float}
    \attr{nights}{Int}
};

% ---------- ROW 3: booking_treatments, payments ----------
\node[entity] (BT) at (0, -5) {
    \textbf{booking\_treatments} \nodepart{second}
    \pk{id}{String}
    \fk{bookingId}{String}
    \attr{name}{String}
    \attr{quantity}{Int}
    \attr{unitPrice}{Float}
    \attr{totalPrice}{Float}
};

\node[entity] (Pay) at (6.5, -5) {
    \textbf{payments} \nodepart{second}
    \pk{id}{String}
    \fk{bookingId}{String}
    \attr{amount}{Float}
    \attr{method}{PaymentMethod}
    \attr{status}{PaymentStatus}
    \attr{paidAt}{DateTime?}
    \attr{createdAt}{DateTime}
};

% ============================================================
%  ХАРИЛЦААНУУД
% ============================================================

% users ── bookings
\draw[rel] (User.south) -- (Booking.north)
    node[mult, midway]{1 : N};

% sanatoriums ── rooms
\draw[rel] (Sana.east) -- (Room.west)
    node[mult, midway]{1 : N};

% bookings ── booking_rooms
\draw[rel] (Booking.east) -- (BR.west)
    node[mult, midway]{1 : N};

% rooms ── booking_rooms
\draw[rel] (Room.south) -- (BR.north)
    node[mult, midway]{1 : N};

% bookings ── booking_treatments
\draw[rel] (Booking.south) -- (BT.north)
    node[mult, midway]{1 : N};

% bookings ── payments
\draw[rel] ($(Booking.south) + (1.5,0)$) -- ++(0,-2) -| (Pay.north)
    node[mult, pos=0.5]{1 : N};

% N:N тэмдэглэл
\node[font=\footnotesize\itshape, text=violet!70!black]
    at (9.7, 5.8) {bookings $\longleftrightarrow$ rooms \ (N : N)};

\end{tikzpicture}
}

\caption{Өгөгдлийн сангийн ER диаграмм}
\label{fig:er-diagram}
\end{figure}

\subsection{Хүснэгтүүдийн тайлбар}

\begin{table}[H]
  \centering
  \caption{Өгөгдлийн сангийн үндсэн хүснэгтүүд}
  \label{tab:db-tables}
  \begin{tabular}{lp{4cm}ll}
    \toprule
    \textbf{Хүснэгтийн нэр} & \textbf{Тодорхойлолт} & \textbf{Гол түлхүүр} & \textbf{Гадаад түлхүүр} \\
    \midrule
    users               & Системийн бүх хэрэглэгч, үүрэг (role) талбарын утгаар ялгана   & id & — \\
    sanatoriums         & Бүртгэлтэй сувиллуудын нэр, байршил, холбоо барих мэдээлэл     & id & — \\
    rooms               & Сувилл бүрийн өрөөний мэдээлэл, үнэ, хүчин чадал               & id & sanatorium\_id \\
    bookings            & Захиалгын бүртгэл, статус, огноо, нийт үнийн дүн               & id & user\_id, room\_id \\
    booking\_treatments & Захиалгад хамаарах эмчилгээний нэр, тоо, үнийн дэлгэрэнгүй   & id & booking\_id \\
    \bottomrule
  \end{tabular}
\end{table}